% Deutsche Parallelfassung zu foundations/notation.tex

\chapter{Notation als Abhängigkeitssprache}
\label{chap:notation}

Die Notation wird in der Reihenfolge eingeführt, in der das Problem von
CAD-Kurve und Vermessungspolyline sie benötigt. Ein Symbol verdient seinen
Platz, indem es eine Eingabe, eine Ausgabe, eine Qualifikation oder eine
Operation bezeichnet. Es weist weder ingenieurtechnische Identität noch
Autorität zu.

\section{Wo gilt die Konstruktion?}

Das Intervall \(I\subseteq\mathbb{R}\) ist die intrinsische longitudinale
Domäne, und \(s\in I\) ist die Bogenlänge. Ein Strich bezeichnet die
Differentiation nach \(s\), sofern keine andere unabhängige Variable angegeben
ist:
\[
f'(s)=\frac{\dd f}{\dd s}(s).
\]
Das Symbol \(\mathcal K\) bezeichnet den zulässigen Raum der
Krümmungsfunktionen. Die Schreibweise \(\kappa\in\mathcal K\) benennt damit
sowohl die konstruktive Eingabe als auch die Regularitätsbedingungen, die das
anwendbare Modell verlangt.

\section{Was ist gegeben und was wird rekonstruiert?}

Die intrinsische Eingabe ist die Krümmung \(\kappa(s)\). Überhöhungsdrehung
\(\phi(s)\) und ihre Rate \(\omega(s)=\phi'(s)\) werden hinzugefügt, wenn der
räumliche Eisenbahnzustand sie erfordert. Richtungswinkel \(\theta(s)\),
ebene Position \(\gamma(s)\) und lokale Richtungen
\(\mathbf t(s),\mathbf n(s),\mathbf b(s)\) werden durch die benannten
Entwicklungs- und Realisierungsbeziehungen abgeleitet. Das lokale Frame ist
\[
F(s)=\bigl(\mathbf t(s),\mathbf n(s),\mathbf b(s)\bigr).
\]

Ein allgemeiner Zustand ist \(x(s)\in\mathcal X\), eine Steuerung ist
\(u(s)\), und eine bewertete oder abgeleitete Größe liegt in \(\mathcal Y\).
Die Namen \(\mathrm{pose2}(s)\) und \(\mathrm{pose3}(s)\) bezeichnen die in
\cref{chap:state-model} definierten Zustandsbegriffe; sie sind keine Synonyme
für eine gespeicherte Kurve oder Polyline.

\section{Welchen Zustand und welche Quelle beschreibt eine Größe?}

Indizes bewahren die Rolle einer Größe:
\[
(\cdot)_{\mathrm{target}},\qquad
(\cdot)_{\mathrm{real}},\qquad
(\cdot)_{\mathrm{meas}}.
\]
Sie unterscheiden einen beabsichtigten Zustand, einen physisch realisierten
Zustand und eine gemessene Beschreibung, die als Beobachtung verwendet wird.
Weitere Indizes können eine Quelle, ein Frame, eine Koordinatendomäne oder
einen Realization Context benennen. Qualifikation verhindert versehentliche
Substitution; sie beweist nicht, dass zwei qualifizierte Größen dasselbe
Engineering Object betreffen.

\section{In welchem Raum wird ein Vergleich berechnet?}

Gleisbezogene Koordinaten \((s,d,h)\) gehören zu
\(\Omega_{\mathrm{track}}\), während ein realisierter Punkt
\(x_{\mathrm w}\) zu \(\Omega_{\mathrm{world}}\) gehört. Die
Track-to-World- und World-to-Track-Operationen lauten
\[
\mathcal P_{\mathrm{tw}}:
\Omega_{\mathrm{track}}\rightarrow\Omega_{\mathrm{world}},
\qquad
\mathcal P_{\mathrm{wt}}:
\Omega_{\mathrm{world}}\rightarrow\Omega_{\mathrm{track}}.
\]
Eine Koordinatentransformation \(\mathcal C\) überführt eine
World-Space-Beschreibung in eine andere. Sie ist keine
World-to-Track-Interpretation und begründet keine Objektidentität.

\section{Welche Operation erzeugt das nächste Artefakt?}

Die wiederkehrenden Operatoren sind Abtastung \(\mathcal S\), Interpolation
\(\mathcal I\), Bewertung \(\mathcal E\), Frame-Konstruktion \(\mathcal F\),
physische Antwort \(\mathcal M\), lokale Anpassung \(\mathcal A\) und
Realisierung \(\mathcal R\). Jede spätere Verwendung muss Domäne, Kodomäne
und anwendbare Annahmen benennen. Insbesondere erzeugt \(\mathcal E\) ein
Bewertungsergebnis, niemals eine Entscheidung.

Symbole, die nur unter einer Erdoberflächenrealisierung benötigt werden,
werden zurückgestellt, bis diese Qualifikation hergestellt ist:
\(\kappa_g\) und \(\kappa_n\) bezeichnen geodätische Krümmung und
Normalkrümmung, \(\mathbf n_{\mathrm E}\) die Oberflächennormale des
Ellipsoids. Ihre Bedeutung hängt von der anwendbaren Fläche und metrischen
Verbindung ab.

Skalare werden kursiv, Vektoren fett, Mengen und Räume groß oder
kalligrafisch und Operatoren dort kalligrafisch gesetzt, wo dies ihre Rolle
sichtbar macht. Diese typografischen Konventionen unterstützen die
Abhängigkeitssprache; sie ersetzen keine ausdrückliche Typisierung.
